\documentclass[UTF8,11pt,fontset=fandol]{ctexart}

\usepackage[a4paper,top=2.0cm,bottom=2.2cm,left=2.4cm,right=2.4cm]{geometry}
\usepackage{amsmath}
\usepackage{booktabs,tabularx,array}
\usepackage{xcolor}
\usepackage{graphicx}
\usepackage{float}
\usepackage{caption}
\usepackage{xurl}
\usepackage{hyperref}
\usepackage{tikz}
\usetikzlibrary{arrows.meta,positioning,fit,backgrounds,calc}
\usepackage{microtype}

% Three roles only: ink for text, one accent, one rule gray.
% Diagram groups are separated by outline, fill tint and line style
% rather than by hue, so the figures survive black-and-white printing.
\definecolor{ink}{HTML}{1A1D23}
\definecolor{accent}{HTML}{1F3A5F}
\definecolor{muted}{HTML}{6B7280}
\definecolor{rulegray}{HTML}{C3C8D0}

\hypersetup{
  colorlinks=true,
  linkcolor=ink,
  citecolor=accent,
  urlcolor=accent,
  pdfauthor={王锦敖},
  pdftitle={面向多智能体基础能力的可验证 RL 数据构造}
}

% CJK convention is indent-only; parskip stays at zero.
\setlength{\parindent}{2em}
\setlength{\parskip}{0pt}
\renewcommand{\arraystretch}{1.25}
\captionsetup{font=small,labelfont=bf}
\setcounter{secnumdepth}{2}
\setcounter{tocdepth}{2}
\ctexset{
  section={
    format=\large\bfseries\color{ink},
    beforeskip=1.5em,
    afterskip=0.6em
  },
  subsection={
    format=\normalsize\bfseries\color{ink},
    beforeskip=1.0em,
    afterskip=0.4em
  }
}
\pagestyle{plain}

\begin{document}

\hypersetup{pageanchor=false}
\thispagestyle{empty}

\vspace*{1.5cm}
\noindent{\bfseries\fontsize{21}{28}\selectfont\color{ink}
基于信息缝与权限划分的\\[0.15em]
可验证多智能体 RL 任务构造\par}

\vspace{0.55cm}
\noindent{\large\color{ink}一种把智能体任务库改造成多智能体 RL 任务的方法，并在 Gaia2 上完整实现\par}

\vspace{0.9cm}
\noindent{\large\color{ink}王锦敖\par}
\vspace{0.35em}
\noindent{\color{muted}Multi-Agent Task Forge\par}
\vspace{0.2em}
\noindent{\color{muted}2026 年 7 月 17 日\par}

\vspace{0.5cm}
\noindent{\color{rulegray}\rule{0.32\textwidth}{0.4pt}}

\vspace{0.45cm}
\noindent\begin{minipage}{0.9\textwidth}
\setlength{\parindent}{0pt}
{\bfseries 摘要}\hspace{0.6em}本文给出的是一套方法，以及一个被完整跑通的实例。该方法把原本并非为多智能体设计的智能体任务库，改造成必须通过协作才能完成的训练任务，同时完全保留原任务的环境与判定器，不新写任务、不新写环境、更不新写奖励。它对任务库只提出四项要求：判定器由他人编写、环境可安装可执行、每条任务附带参考信息（ground truth）、动作面可以沿角色边界切开。任何满足这四项的任务库都可以套用本方法；Gaia2——Meta 的 Agents Research Environments 所附带的基准\cite{meta2025are}——是本文实际运行的那一个，而 \S\ref{sec:substrate} 会明确指出：它也是唯一的一个。具体到 Gaia2：每个场景（scenario）都附带一组黄金写操作（gold write actions）和完整的应用初始状态；我们据此做\emph{状态溯源}判定——某个写操作的参数值若出现在另一个应用的初始状态中、不在目标应用自身状态中、也没有出现在用户指令里，那么这条信息必然跨过了应用边界。每道题需要哪些角色，完全由参考信息实际触及的应用决定。改造后，主智能体不再持有任何应用工具，只能把工作委派给各应用的专属智能体（specialist）；参考信息和判定器则保存在主智能体无法访问的辅助容器（sidecar）中。最终判定仍由 Gaia2 原有的写操作判定器给出，其脚本模式不含任何大模型。

\vspace{0.5em}
这条链路中，批量生成的开销并不高：挖掘、角色划分和 Harbor 目录渲染都由静态分析完成，37 个入选场景可以直接生成 148 个配置实例。真正的成本发生在生成之后，而本文的主要内容正是这笔成本买到了什么。\S7 报告了覆盖全部 148 个实例的单种子实验，其结论是否定的：无约束对照组只在 37 个场景中的 2 个上成功，各约束组的均值与对照组相差不超过 $0.035$，因此 \S4.3 的接受判据\textbf{拒绝给出任何结论}——当对照组本身没有越过下界时，任何叠加其上的约束都无法定价。所以下一步不是加旋钮，也不是加实例，而是先让对照组跑通：这批实验中有 49\% 的轨迹终止于格式错误的工具调用或环境提前结束，协作能力根本没有进入被考察的范围。
\end{minipage}

\clearpage
\thispagestyle{empty}
{\setlength{\parskip}{0pt}
\renewcommand{\contentsname}{\normalsize\bfseries 目录}
\tableofcontents}

\clearpage
\hypersetup{pageanchor=true}
\pagenumbering{arabic}

\section{问题与目标}

一条强化学习任务通常由指令、可交互环境和奖励三部分组成。对单智能体任务来说，验证往往很直接：文件是否生成、数据库是否更新、测试是否通过，都可以由程序检查。多智能体任务多了一层要求。除了完成最终目标，主智能体还要找到合适的分工，并把必要的信息准确地传给不同角色。

角色数量并不能自动带来协作。如果主智能体仍能直接调用全部工具，或任意一个子智能体都能独立完成任务，那么多出来的角色只是界面变化，并没有改变问题本身。另一种做法是直接评价对话是否“像在合作”，但这种奖励高度依赖主观判断，也很难稳定地扩展到大规模数据。

我们的选择是从已有可靠判定器的真实任务出发，只改变工具权限和信息分布，使协作成为完成任务的必要条件。这样既能迫使模型规划、委派和传递信息，又能继续用最终环境状态计算奖励。

把这个目标提到它真正应当被复用的层次上，它就不是一批数据，而是一套\textbf{方法}。目前并不存在“多智能体轨迹 + 环境 + 免费判定器”的现成语料，因此凡是要造这类数据的人，都必须自己制造多智能体结构；真正的选择只有一个：\emph{制造哪一层}。制造判定器，错误就会隐形——因为你自己的测试会认同你自己的错误，\S4.1 正是这一失败被测出来的记录。只制造\emph{约束}——它改变谁能做什么，从不改变什么算做完——失败模式就退化成“题目太易或太难”：可见，且可以用 \S4.3 的判据定价。本文余下的部分，就是这套方法的四个动作——准入任务库（\S6.6）、挖掘信息缝（\S3.1）、施加约束（\S3.2）、只接受落在下界与对照之间的实例（\S4.3）——以及被这四个动作完整跑通的一个任务库：Gaia2。

\section{能力选择与学习顺序}

多智能体协作涉及多种能力，但并不是每一种都适合直接作为 RL 维度。为了得到清晰的训练信号，我们优先选择能够由环境结果验证、并且可以逐层组合的能力。MAST 将常见失败归纳为系统设计、智能体间失配、任务验证和终止等类别\cite{cemri2025mast}；MultiAgentBench 则比较了星形、链式、树形和图结构等通信拓扑下的协作表现\cite{zhu2025multiagentbench}。这些因素彼此相关，造题时需要分阶段引入，而不是一次全部叠加。

\begin{table}[H]
\centering
\small
\caption{当前优先训练的能力}
\label{tab:capabilities}
\begin{tabularx}{\textwidth}{>{\raggedright\arraybackslash}p{2.65cm} X >{\raggedright\arraybackslash}p{4.05cm}}
\toprule
能力 & 模型需要学会什么 & 如何观察失败\\
\midrule
任务拆解与依赖识别 & 把一个目标拆成若干可执行子任务，并判断哪些中间结果必须先获得 & 遗漏必要应用、顺序错误、下游缺少上游信息\\
角色分配 & 根据子智能体拥有的信息和工具，把子任务交给合适的角色 & 请求发给错误角色，或要求一个角色访问它看不到的信息\\
跨角色通信 & 在委派中补充实体、时间范围和上一步结果，使局部执行服务于全局目标 & 子智能体完成了自己的请求，但最终世界状态仍然不正确\\
并行调度 & 在满足依赖和资源约束的同时减少不必要的串行等待 & 资源冲突、漏排任务、依赖未满足或总完成时间过长\\
\bottomrule
\end{tabularx}
\end{table}

当前交付直接覆盖表~\ref{tab:capabilities} 中的前三项。并行调度虽然也可以通过环境状态验证，却不在本次交付范围内：主智能体每轮只能发出一次 \texttt{DELEGATE}，\texttt{run\_specialist} 返回之后才能发出下一次，因此无论场景的依赖图有多宽，委派始终是串行阻塞的。Gaia2 至少让缺失的另一半变得容易设想：它的时钟是模拟的，主智能体用 \texttt{WAIT} 显式推进时间，因此总完成时间本身就是环境已经在记录的量。真正缺的是一个能同时发出两条委派的运行时，以及一个为这种差异定价的奖励项。在两者都具备之前，这一维度暂不纳入。

课程可以按依赖关系逐步展开。第一层只要求模型识别角色并完成一次正确委派；第二层加入跨应用依赖，要求主智能体把上游结果带到下游；第三层隐藏角色的能力说明，让主智能体自行判断“谁可能知道什么”；之后再加入委派预算、更多角色和更复杂的通信拓扑。第三层实际上已经涉及一部分心智理论（Theory of Mind）：主智能体必须估计其他角色掌握的信息，而不能默认对方与自己共享上下文。不过，当前奖励只能判断最终任务是否完成，无法确认一次失败究竟来自心智建模、任务拆解还是消息表达，因此这里不把心智理论和动态恢复单独列成训练维度。

更细的失败归因仍然可以作为诊断工具。例如，让教师模型阅读交互记录，判断问题出在拆解、路由还是沟通。但这相当于新增了一个判定器，必须先与人工标注对照，报告一致率，才能知道它有多可信。在完成这项验证之前，终态分数仍应是唯一的训练奖励，教师模型最多用于分析失败和设计后续课程。本文尚未进行这项研究，因此表中只保留能够由终态直接验证的能力。

\begin{quote}
\small
\noindent 建议的进阶顺序为：单次角色路由 $\rightarrow$ 星形拓扑、能力文档可见 $\rightarrow$
星形拓扑、仅角色名称可见 $\rightarrow$ 有限预算与多级依赖 $\rightarrow$ 链式/树形拓扑与运行时故障。
\end{quote}

当前实现覆盖到第三层。委派预算、多级依赖和复杂拓扑已有设计方向，但在进入训练集之前，还需要分别验证可解性和稳定的难度增量。

\section{从原始任务到协作任务}

核心思路可以概括为：\emph{保留任务和判定器，只重构智能体访问任务的方式。}

原始任务已经给出了用户指令、交互环境和程序化判定器。改造时，我们只在环境外层加入权限、可见性、通信拓扑和委派预算等约束。模型可以采用不同的拆解和通信策略，但最终仍由原任务的测试判断是否完成。

更具体地，记原始任务为
\[
  \tau=(x,\mathcal{E},J),
  \qquad
  \mathcal{E}_c=W(\mathcal{E};c),
  \qquad
  \tau_c=(x,\mathcal{E}_c,J).
\]
其中 $x$ 表示用户指令，$\mathcal{E}$ 表示原始环境，$J$ 是程序化判定器。包装算子 $W$ \emph{只改变环境}，不改写指令和判定器，因此等式两边的 $x$ 与 $J$ 保持不变。约束
\[
  c=(P,V,G,B)
\]
的四个分量是：$P$ 权限划分（哪个角色持有哪个应用的 API），$V$ 角色可见性（主智能体能够看到哪些角色信息），$G$ 通信图（谁可以向谁发消息），$B$ 委派预算。

这里的 $P$ 描述划分方式：主智能体不持有任何应用工具，每个专属智能体只绑定一个应用。具体需要哪些应用，不由配置预先指定，而是从参考信息中推导（\S3.1）。权限划分是其余约束成立的前提；如果主智能体仍能直接访问全部工具，可见性、拓扑和预算都不会真正限制它。判定器 $J$ 读取辅助容器中的终态，而 $W$ 保证主智能体所在的文件系统中没有 $J$ 或完整参考信息（\S5.1）。

这一思路与 SWE-smith 复用仓库既有测试的原则相近\cite{yang2025swesmith}：优先继承可靠的验证逻辑，而不是为每道新题重新定义“什么算完成”。整个过程分为三步：识别任务中的跨应用依赖，据此确定角色集合，再把同一任务渲染成不同的协作配置。

\subsection{从真实依赖中选择题目}

Gaia2 提供一组日常应用——Calendar、Contacts、InternalContacts、Messages、Emails、Chats、Cabs、RentAFlat、Shopping 等——它们运行在一个时钟连续流动、事件异步触发的模拟世界中\cite{meta2025are}。它通过比对智能体的写操作与场景的黄金写操作来判断任务是否完成，并提供一个不含大模型的脚本判定模式。本项目读取两个公开划分：\texttt{mini} 与 \texttt{adaptability}，各 160 个场景。

这里的参考信息不是代码，因此两条挖掘规则需要换一套机制来实现。角色集合仍然完全由任务本身决定：它是黄金写操作触及的应用，\emph{再加上}被信息缝证明确实被读取过的应用。第二个条件是数据本身逼出来的——Gaia2 的判定事件只记录写操作，一个仅被读取的应用不会在参考信息中留下任何痕迹，否则就会被漏掉。用于向用户汇报的 \texttt{AgentUserInterface} 不提供业务信息，因此不计入角色：它是本任务库的提交通道，\S4.4 会给出把提交通道当作协作者会让数字变成什么样。

按这一规则，\texttt{mini} 的 160 个场景中有 111 个、\texttt{adaptability} 的 160 个中有 158 个确实涉及至少两个应用；其余直接过滤，不额外添加无关角色。

\subsection{把单智能体任务改成协作任务}

改造后的环境包含一个主智能体（Main）和若干应用专属智能体。主智能体能看到完整用户目标，却不能调用任何应用工具；Contacts 专属智能体只能访问 Contacts，Cabs 专属智能体只能访问 Cabs。每个专属智能体只接收主智能体发来的局部请求，并不知道完整任务。

以入选场景 \texttt{scenario\_universe\_25\_b932pi} 为例，它的角色集合有 5 个应用，实测依赖图深度为 3。用户的要求是：为今天最早一场日程的参会人订一辆默认档位的车，发车时间比日程早 10 分钟，\emph{上车地点是这位参会人的住址}；下车地点取决于该日程是否带 work 标签；订好后用 Messages 通知参会人；随后监控这笔订单 4 分钟，若司机取消，则删除该日程并以不含参会人的方式重建，再次通知各方。

上面被强调的那一句就是信息缝，值得把理由说清楚。黄金写操作 \texttt{Cabs.order\_ride} 携带的参数是 \texttt{start\_location = "Sint-Michielsstraat 6, 3000 Leuven, Belgium"}。这个字符串既不在用户指令里（用户只说了“他的住址”），也不在 Cabs 自身的初始状态里，而是存在于 Contacts 中。也就是说，这条信息\emph{必须}跨过应用边界，这次写操作才可能正确——而这正是状态溯源判定所检测的，整个推断过程无需阅读场景的自然语言描述。信息传递路径由任务本身决定：
\[
  \text{Calendar 给出最早日程、标签与参会人}
  \longrightarrow
  \text{Contacts 给出该参会人的住址}
\]
\[
  \longrightarrow
  \text{主智能体应用 work 标签的下车规则}
  \longrightarrow
  \text{Cabs 下单}
  \longrightarrow
  \text{Messages 通知}.
\]

任何一个专属智能体都无法独立完成：Calendar 知道时间和人，却不能订车；Contacts 知道住址，却没有可下单的工具；Cabs 能下单，却既不知道时间也不知道地址。Gaia2 还引入了静态世界无法表达的一个维度——场景的后半段取决于 4 分钟后是否发生取消，因此主智能体必须停止动作并主动推进时钟（\S5.1），而不是靠猜。

这类任务同时要求模型判断信息所在的角色、安排有依赖的执行顺序，并在委派时带上足够的上下文。题面仍是一条自然语言指令，依赖结构只存在于环境和参考信息中，不会直接展示给模型。需要注意的是，终态奖励只能判断整体结果，无法把一次失败单独归因到上述某项能力（见 \S8.1）。

\section{造题方法：四个关键环节}

SWE-smith 最值得借鉴的并不是如何注入 bug，而是它把验证问题交还给仓库原有的测试：先破坏一个本来能运行的仓库，再用既有测试判断修复是否成立\cite{yang2025swesmith}。沿着这条思路，一套可规模化的造题方法需要回答四个问题：什么保持不变、候选如何生成、什么样的候选可以保留，以及每种生成方式最终能留下多少有效任务。SWE-smith 对五类算子分别测得 56\%、35\%、40.2\%、33.8\% 和 96.9\% 的良率。

下面按照这四个环节说明本项目的设计，也明确区分已经实现的部分和仍需验证的部分。

\subsection{不变量：优先继承已有判定器}

造题既可以重建任务世界和判定器，也可以只改变智能体受到的约束。两种做法的风险并不相同。约束设置不当，任务会变得过易或不可解，这通常能从对照实验中发现；判定器一旦写错，奖励就会稳定地衡量错误目标，而且自测未必能暴露问题。因此，本项目只改变“谁能访问什么”，不重新定义“什么算完成”。这样做的目的不是简化实现，而是尽量把失败留在可观测、可比较的部分。

这并不意味着语言模型不能出现在造题链路中。它可以生成候选任务，也可以充当环境中的专属智能体，或在没有程序化判定器时辅助诊断。区别在于：一旦模型参与打分，就需要先用人工标注测量其一致率。这个要求同样适用于任何自建规则，包括 \S4.4 用来筛选任务的信息缝分析；它目前已有代码和结果，但误差率还没有单独评估。

因此，已有的程序化判定器仍是首选。其他方案并非不可行，只是需要额外承担判定器验证的成本。

\paragraph{源环境选择。}
源环境需要满足几个基本条件：有可复用的程序化判定器；环境真实、可安装、可执行；至少训练集提供完整参考信息，以便分析任务结构和验证可解性；工具空间还要能够沿两个以上的角色边界切分。

实践中还需要再加一条：角色边界必须同时形成信息边界。若两个角色可以各自完成工作，只需安排先后顺序，那么任务考查的是分派，而不是信息协作。只有当下游动作依赖上游角色独有的事实时，委派和通信才会影响最终结果。本文把这种跨角色的信息依赖称为“信息缝”。共享同一文件系统的软件仓库通常较难形成这种边界，因为信息对所有角色同时可见，角色划分容易退化为文件分工。

\subsection{算子：如何产生候选任务}

候选任务的生成方式分为程序化变换、语言模型改写和任务组合三类，见表~\ref{tab:operators}。

\begin{table}[H]
\centering
\small
\caption{候选生成方式与当前实现状态}
\label{tab:operators}
\begin{tabularx}{\textwidth}{>{\raggedright\arraybackslash}p{1.6cm} >{\raggedright\arraybackslash}p{2.2cm} X >{\raggedright\arraybackslash}p{1.5cm} >{\raggedright\arraybackslash}p{1.5cm} >{\raggedright\arraybackslash}p{1.3cm}}
\toprule
产生方式 & SWE-smith 的对应 & 本项目的算子 & 改题/改配置 & 良率 & 状态\\
\midrule
程序化 & Procedural（13 种 AST 变换） & \texttt{partition}：按参考信息触及的应用划分角色，主智能体不持有任何工具 & 改配置 & 入选 \textbf{37} 场景；实例无法判定 & 已交付\\
程序化 & Procedural & \texttt{blind}：隐藏角色能力说明，只留角色名 & 改配置 & 实例无法判定 & 已交付\\
程序化 & Procedural & \texttt{route}（委派经济）：按任务推导出的软性 \texttt{b}$N$ 委派目标 & 改配置 & 实例无法判定 & 已交付\\
程序化 & Procedural & \texttt{route}（链式/树形拓扑）：改变角色之间的通信路径 & 改配置 & --- & 未实现\\
程序化 & Procedural & \texttt{perturb}：运行中撤下一个角色，迫使重规划 & 改题 & --- & 未实现\\
语言模型改写 & LM Modify / LM Rewrite & \texttt{weaken}：抽走指令中的日期与实体，移入辅助容器的隐藏用户状态 & 改题 & --- & 未实现\\
组合 & Combine & \texttt{combine}：把共享实体的两道题合成一道 & 改题，判定器需重构 & --- & 未实现\\
\bottomrule
\end{tabularx}
\end{table}

“已交付”表示该配置出现在已发布的实验中；“未实现”则表示代码链路还不完整。表中的“实例无法判定”不等于良率为零，也不应被这样解读：它的含义是 \S4.3 的接受判据\emph{拒绝给出结论}，因为本任务库上的无约束对照自身尚未越过下界（\S7）。对照失败时，旋钮只是未被定价，而不是被证明无效。链式拓扑仍未实现：专属智能体之间的交接关系虽有定义，却没有实际调用点，因此该拓扑下的场景不可解，其难度信号是假的。表中的算子名称用于描述方法，对应代码中的 \texttt{Topology}、\texttt{Visibility} 和软性 \texttt{b}$N$ 目标等参数。

需要区分“新任务”和“新配置”。\texttt{partition}、\texttt{blind} 与委派经济目标只改变同一个场景的访问方式，因此 37 个场景乘以 4 个实例得到的是 148 个运行单元，而不是 148 个彼此独立的新题。SWE-smith 的 50k 条实例来自对任务本身的修改。若要在本项目中获得同样意义上的题量增长，就必须引入会改变任务内容的算子，并重新测量其良率。

\texttt{combine} 尤其需要谨慎。Gaia2 是拿场景的写操作去比对\emph{该场景}的黄金动作集，因此合并两个场景就必须重新决定哪些写操作是合法的，简单地取 $J_A\wedge J_B$ 并不能得到自洽的新判定器。组合任务可能带来更高产量，也最容易破坏“继承判定器”这一前提。

\subsection{判据：什么样的候选可以保留}

SWE-smith 的保留规则很直接：应用候选补丁后运行测试，只有原本通过的测试变红，候选才算有效。整个过程由程序完成，不需要额外判断。

本项目使用类似的夹逼规则。记无动作基线为 $r_0$，无约束对照为 $r_{\mathrm{open}}$，某个受约束配置的得分为 $r$。只有满足
\[
  r_0 < r < r_{\mathrm{open}}
\]
并且 $r_{\mathrm{open}}>r_0$ 时，这个配置才提供可解释的训练信号。若 $r\leq r_0$，它在奖励上无法与不采取行动区分，原因可能是任务过难，也可能是运行链路故障；若 $r\geq r_{\mathrm{open}}$，说明约束没有带来可测影响，得到的仍是基础任务能力；若无约束对照本身落在基线上，则基础任务或运行环境已经有问题，不能再从受约束分数中解释难度。

可解性需要单独验证，不包含在上述夹逼规则中。若任务库的参考信息是代码，可解性可以靠实际运行参考解来证明。Gaia2 给出的是黄金\emph{写操作}而不是可执行脚本，因此可解性只能由无约束对照代为证明：一个 \texttt{open} 实例若能成功，就说明该场景对持有全部工具的智能体是可达的。这个替代方案比参考解弱，本文并不掩饰这一点——它正是 \S7 中否定结论的直接成因：当对照失败时，可解性与难度就无法区分了。

这条规则已经写入造题链路。项目早期曾把一个下界分数解释为权限划分带来的难度，后来才确认它只是无动作基线。它当下最有用的性质恰恰在后文得到体现：与其给出一个看起来像结论的数字，它选择拒绝判定。

\subsection{良率：从候选到可用数据}

良率描述候选中最终有多少可以使用。本项目有两次筛选，因此需要分别看两个比例：一是题库中有多少任务具备真实的跨角色信息依赖，二是这些任务在具体配置下有多少能产生训练信号。

\subsubsection*{选题良率：多应用不等于可协作}

宽松的筛选条件是 \texttt{len(roster) >= 2}。它只能说明参考信息触及了至少两个应用，却不能证明两个应用之间存在信息依赖。若第二个应用只是独立执行另一项操作，增加角色并不会真正考查通信。

在 Gaia2 上，这一点可以用状态溯源来测量，且不需要任何模型调用：取黄金写操作中的每一个参数值，若某个值（一）出现在\emph{另一个}应用的初始状态中，（二）不在目标应用自身状态中，（三）也没有出现在用户指令里，则保留该场景。三个条件缺一不可——没有（二），这个值本来就可能是本地的；没有（三），用户已经直接告诉了智能体，信息并不需要跨越边界。

\begin{table}[H]
\centering
\small
\caption{两个 Gaia2 划分上的选题良率，各 160 个场景（\texttt{sweep/gaia2\_mini\_admission.json}、\texttt{sweep/gaia2\_adaptability\_admission.json}）}
\label{tab:selection}
\begin{tabularx}{\textwidth}{X c c}
\toprule
过滤 & \texttt{mini} & \texttt{adaptability}\\
\midrule
宽松筛选：把用户接口通道也算作协作者 & 160/160（100\%） & 160/160（100\%）\\
严格筛选：剔除提交通道，要求角色集合 $\geq2$ & 111/160（69\%） & 158/160（99\%）\\
\textbf{+ 信息缝}：要求事实真的跨应用流动 & \textbf{17/160（11\%）} & \textbf{30/160（19\%）}\\
\bottomrule
\end{tabularx}
\end{table}

两轮过滤解决的是同一类问题，只差一个层级：第一轮排除“不承载信息的提交通道”，第二轮排除“不承载信息的第二个应用”。在这里，第二轮的代价要大得多。\texttt{adaptability} 上 99\% 的场景是多应用的，而只有 19\% 存在信息缝——也就是说，\textbf{看起来可协作的场景里有四分之三是装饰性的}；一条止步于角色数检查的链路，会把它们全部当成训练数据。

两个划分有 10 个场景重叠，因此语料共包含 \textbf{37 个不重复的信息缝场景}，渲染为 148 个实例，角色集合 2--7 个应用，黄金写操作 5--16 条，依赖图深度 1--4，并行宽度最高 12。\texttt{mini} 各类别的分布为：adaptability 10/32，time 4/32，execution 3/32，ambiguity 0/32，search 0/32。后两行为零属于预期而非异常——search 与 ambiguity 场景以读取为主，而 Gaia2 只验证写操作，一个最终落在\emph{回答}里而不是落在写操作里的事实，对基于写参数的溯源判定是不可见的。

关于这个数字有两点必须说明。其一，该启发式\textbf{在构造上就会低估}：它做的是精确字符串匹配，因此中途被改写过格式的日期不可见，长度不足 8 的值因易碰撞而被忽略，出现在三个及以上应用中的值被视为环境噪声。其二，它是本项目自行设计的过滤器，按 \S4.1 的规则，它同样欠一份与人工标注对照的验证研究。这项研究尚未进行。目前支持它的只有结构性证据：最密集的缝边是 Contacts$\rightarrow$Cabs、InternalContacts$\rightarrow$Cabs 与 RentAFlat$\rightarrow$Cabs，这更像是应用本身的性质而不是随机散布——但“看起来合理”不是测量。

\subsubsection*{配置良率：目前还无法测量，而这正是结论}

把 \S4.3 的夹逼规则应用到 \S7 的实验上，得到的结果是\emph{没有结果}，原因只有一个：$r_{\mathrm{open}}$ 没有越过下界。无约束对照在 37 个场景中只成功 2 个，平均部分分为 $0.251$，而各约束组与它的差距都不超过 $0.035$。

判据的第三个条件正是为这种情形准备的。把 \texttt{star-docs} 的 $0.253$ 与对照的 $0.251$ 相比并宣称“权限划分是免费的”，与本项目早年把下界当作难度发布，是同一个错误的两个方向。对照失败时，受约束分数不携带关于约束的任何信息，因为没有任何证据表明该场景原本是可达的。

所以本任务库上的配置良率是\textbf{未定义}，而不是零，这个区别决定了下一步该做什么。良率为零意味着旋钮产出的数据不可用；未定义则意味着实验还没有真正开始——评测链路、运行时或基础难度在旋钮受审之前就消耗掉了这些轨迹，而 \S7 会指出是哪一个。它唯一排除掉的，也正是唯一值得排除的：把表~\ref{tab:results} 中的任何一个均值当作已定价的约束来报告。

\section{运行环境、轨迹与验证}

\subsection{容器边界}

Harbor 把每道题组织成独立任务，目录中包含指令、容器环境和测试\cite{harbor2026tasks}。主智能体所在的容器只安装轻量的 \texttt{team} 客户端；Gaia2 世界、专属智能体和任务状态全部运行在辅助容器中。主智能体因而只能通过委派改变环境，无法直接调用应用。

渲染出的 Gaia2 任务有一个性质值得单独说明，因为相对于“世界必须先安装”的任务库这是实实在在的改进：任务把整个世界以 \texttt{environment/scenario.json} 的形式随身携带。运行时既不需要下载数据集，也不需要安装步骤，因此一个任务目录可以完全依靠自身内容复现。

主智能体的动作集里有一个由本任务库逼出来的新增项。Gaia2 的时钟连续流动、事件异步触发，因此一个场景的后半段可以取决于尚未发生的事情。为此主智能体持有 \texttt{WAIT <秒数>}，它推进模拟时间并投递这段时间内到达的通知与用户消息。\emph{何时停止动作}是一个协作决策，因此它属于主智能体而不属于评测框架——若没有它，面对 \S3.2 那个场景的智能体，对四分钟后的触发条件只能靠猜。

\begin{figure}[H]
\centering
\begin{tikzpicture}[
  font=\footnotesize,
  component/.style={draw=rulegray,rounded corners=2pt,fill=white,
    minimum height=0.9cm,text width=2.7cm,align=center,inner sep=5pt},
  main/.style={component,draw=accent,thick,fill=accent!6},
  side/.style={component,draw=ink,thick,fill=white},
  world/.style={component,draw=ink,thick,fill=rulegray!35},
  verify/.style={component,draw=muted,thick,dashed,fill=white},
  boundary/.style={draw=rulegray,rounded corners=5pt,thick,inner sep=11pt},
  edge/.style={-{Latex[length=2.2mm]},semithick,draw=ink},
  verifyedge/.style={-{Latex[length=2.2mm]},semithick,dashed,draw=muted},
  flowlabel/.style={font=\scriptsize,fill=white,inner sep=2pt,text=ink},
  verifylabel/.style={flowlabel,text=muted}
]
  \node[main] (agent) at (0,2.4) {主智能体 Main\\看到完整任务\\无应用工具};
  \node[main] (team) at (0,0) {\texttt{team} 客户端\\\texttt{roster/docs}\\\texttt{ask/wait/done}};
  \node[verify] (verifier) at (0,-3.3) {\texttt{tests/verify.py}\\验证阶段才注入\\token 仅此处可见};

  \node[side] (gateway) at (4.8,0) {会话网关\\执行角色、拓扑\\与委派经济目标\\保存交互记录};
  \node[side] (contacts) at (8.3,1.05) {Contacts 专属智能体\\仅可访问 Contacts};
  \node[side] (cabs) at (8.3,-1.05) {Cabs 专属智能体\\仅可访问 Cabs};
  \node[world] (gaia) at (12.0,0) {Gaia2 世界\\应用工具、模拟时钟\\异步事件};
  \node[world] (judge) at (12.0,-3.3) {写操作判定器\\黄金动作对比观测动作\\返回 pass/failure};

  \node[boundary,fit=(agent)(team)(verifier),
    label={[font=\small\bfseries,text=accent]above:主容器}] (mainbox) {};
  \node[boundary,fit=(gateway)(contacts)(cabs)(gaia)(judge),
    label={[font=\small\bfseries,text=ink]above:maf-env 辅助容器}] (sidebox) {};

  \draw[edge] (agent) -- (team);
  \draw[edge] (team.east) -- node[flowlabel,above=2pt]{\texttt{ask / wait / done}} (gateway.west);
  \draw[edge] (gateway.east) to[out=0,in=180] (contacts.west);
  \draw[edge] (gateway.east) to[out=0,in=180] (cabs.west);
  \draw[edge] (contacts.east) to[out=0,in=180] (gaia.west);
  \draw[edge] (cabs.east) to[out=0,in=180] (gaia.west);
  \draw[edge] (gaia.south) -- node[flowlabel,right=2pt]{终态检查} (judge.north);

  \draw[verifyedge] (verifier.east) to[out=15,in=-115]
    node[verifylabel,pos=0.55,above=2pt]{\texttt{GET /state + token}} (gateway.south);
  \draw[verifyedge] (judge.west) --
    node[verifylabel,below=2pt]{奖励与分项结果} (verifier.east);
\end{tikzpicture}
\caption{Harbor 任务的容器架构。左侧是主智能体可见的容器；Gaia2 世界、专属智能体、交互记录和判定逻辑均位于辅助容器。图中画了两个专属智能体，实际挖掘出的角色集合最多可达七个。虚线表示只在任务结束后启动的验证流程。}
\label{fig:architecture}
\end{figure}

图~\ref{fig:architecture} 中有两条关键隔离边界。主智能体的文件系统不包含世界、黄金写操作或判定器，只能通过 \texttt{team} 发送自然语言请求。验证令牌在任务执行期间同样不可见；Harbor 只在任务结束后注入测试，由验证器读取辅助容器中的终态。权限限制由容器环境强制执行，主智能体无法在执行过程中读取自己的得分。

还有一个与边界有关的工程细节，因为它影响了交付方式：Gaia2 的官方评测框架依赖 pydantic~2，与本项目其余部分的版本固定冲突，因此所有触及 \texttt{are.simulation} 的代码都运行在独立的虚拟环境中，而静态分析——挖掘、准入、渲染——仍在主环境中运行。这与辅助容器是同一个思路，只不过施加在依赖图而不是文件系统上。

\subsection{一条任务如何运行}

图~\ref{fig:trajectory} 给出了 \S3.2 中那个场景 \texttt{scenario\_universe\_25\_b932pi} 的黄金轨迹。一次 \texttt{team ask} 内部可能包含专属智能体的多轮推理和工具调用，图中只展示角色之间传递的信息。任务能否完成，取决于主智能体是否在下单之前先取得参会人的住址，再把住址、发车时间和已经解算好的下车规则一并交给 Cabs；消息数量本身并不重要。

这是该场景的五条黄金写操作所描述的轨迹。必须直说的是，\S7 的实验并没有跑出它：在这个场景上，四个实例——\emph{包括无约束对照}——都只覆盖了 5 条黄金写操作中的 2 条，全部判定为失败。因此这张图是一份规格说明，而不是一条实际轨迹。

\begin{figure}[H]
\centering
\begin{tikzpicture}[
  font=\footnotesize,
  actor/.style={draw=ink,rounded corners=2pt,minimum width=2.1cm,
    minimum height=0.7cm,align=center,font=\small\bfseries,fill=white,
    text width=1.9cm,inner sep=3pt},
  lifeline/.style={draw=rulegray,densely dashed},
  msg/.style={-{Latex[length=2mm]},semithick,draw=ink},
  msglabel/.style={font=\scriptsize,fill=white,inner sep=2pt,align=center,
    text=ink,text width=3.2cm},
  state/.style={draw=ink,rounded corners=2pt,fill=rulegray!35,
    text width=2.2cm,align=center,font=\scriptsize,inner sep=4pt},
  phase/.style={font=\small\bfseries,text=ink}
]
  \node[align=left,anchor=west,font=\scriptsize,text width=10cm] (task) at (-1.05,1.5)
    {用户任务\quad 为今天最早一场日程的参会人订一辆默认档位的车，比日程早 10 分钟发车，上车地点是这位参会人的住址；下车地点取决于日程标签。订好后发消息确认，并持续监控该订单。};

  \node[actor,draw=accent,fill=accent!6] (main) at (0,0) {Main};
  \node[actor] (calendar) at (3.5,0) {Calendar\\专属智能体};
  \node[actor] (contacts) at (7.0,0) {Contacts\\专属智能体};
  \node[actor] (cabs) at (10.5,0) {Cabs\\专属智能体};

  \foreach \x in {0,3.5,7.0,10.5}
    \draw[lifeline] (\x,-0.5) -- (\x,-13.9);

  \node[phase,anchor=west] at (-1.05,-0.9) {执行阶段};

  \draw[msg] (0,-2.2) -- node[msglabel,above=2pt]
    {\texttt{team ask Calendar}\\今天最早的日程是哪一场，标签是什么，谁参加？} (3.5,-2.2);
  \draw[msg] (3.5,-3.4) -- node[msglabel,above=2pt]
    {09:00，带 Work 标签，一位参会人} (0,-3.4);
  \node[state,anchor=west] at (11.85,-3.4)
    {$E_0$：日程、联系人与订单均未修改};

  \draw[msg] (0,-4.6) -- node[msglabel,above=2pt]
    {\texttt{team ask Contacts}\\这位参会人的住址是什么？} (7.0,-4.6);
  \draw[msg] (7.0,-5.8) -- node[msglabel,above=2pt]
    {\textbf{Sint-Michielsstraat 6}——信息缝：用户从未说出这个地址} (0,-5.8);

  \draw[msg] (0,-7.2) -- node[msglabel,above=2pt]
    {\texttt{team ask Cabs}\\携带住址、08:50 与 work 标签下车规则} (10.5,-7.2);
  \draw[msg] (10.5,-8.5) -- node[msglabel,above=2pt]
    {订单已创建} (0,-8.5);
  \node[state,anchor=west] at (11.85,-8.5)
    {$E_1$：订单已下，上车点取自 Contacts};

  \node[state,anchor=west,text width=3.0cm] at (0.35,-9.9)
    {\texttt{team wait 240}：推进时钟，观察是否被取消};

  \draw[msg] (0,-11.6) -- node[msglabel,above=2pt]
    {\texttt{team done completed}} (10.5,-11.6);

  \draw[rulegray,semithick] (-1.05,-12.6) -- (14.2,-12.6);
  \node[phase,anchor=west] at (-1.05,-13.0) {验证阶段};
  \draw[msg,dashed,draw=muted] (10.5,-13.7) -- (0,-13.7);
  \node[msglabel,text width=6.5cm,anchor=south] at (5.25,-13.6)
    {写操作与 5 条黄金动作比对 $\Rightarrow$ pass/fail};
\end{tikzpicture}
\caption{黄金轨迹中的信息流。上游角色取得事实，主智能体将其整理成下游可以执行的请求；奖励由到达终态 $E_1$ 的写操作决定，不依赖轨迹文本。图中省略了 Messages 通知与条件性改期的收尾部分；角色集合共五个应用，此处只画出承载信息缝的三个。虚线表示任务结束后的验证阶段。}
\label{fig:trajectory}
\end{figure}

\subsection{奖励与有效性参照}

验证器不评价对话写得是否合理，也不依赖模型自行报告的工作状态。Gaia2 自身的判定是二元的：把智能体的观测写操作与场景的黄金写操作比对，通过或不通过。模型可以采用不同的拆解和通信路径；只要写操作落到位，判定结果就相同。

二元判定作为训练信号过于稀疏，因此本项目在它之外另算了一个分级信号——而这一步恰恰需要接受 \S4.1 对“自行设计之物”的那种审视。记 $G$ 为黄金写操作总数，$c$ 为智能体确实执行了的条数，$e$ 为其中参数与黄金完全一致的条数，则
\[
  r=\frac{e+\tfrac{1}{2}(c-e)}{G}.
\]
完全匹配计满分，工具正确而参数有误计半分：一个把正确消息发给了错误号码的智能体确实完成了部分工作，把它记为零就又回到了全有全无的信号。有两点使它保持诚实。它是\textbf{确定性的、不含模型的}——完全由磁盘上的轨迹重算，任何判定器都无法改变它；它同时是\textbf{官方判定之下的一个下界，而不是它的替代品}——它绝不会认可一个判定器判为失败的场景。但它终究是本项目自己设计的奖励，因此只作为诊断量报告，绝不作为接受判据。\S7 中的每一个判定结论都来自 Gaia2 自己的验证器。

每个配置都需要三个参照。其一用于证明任务本身可解；无约束对照让单个智能体直接访问全部应用；无动作基线则测量在不修改环境时能够自然得到多少分。只有当受约束配置稳定地落在无动作基线与无约束对照之间，它才可能为 RL 提供有效梯度。

在本任务库上，缺的正是第一个和第三个参照，\S7 就是这笔欠账的代价。Gaia2 给出的是黄金写操作而不是可运行的参考解，因此无约束对照被迫兼任可解性证明——而当对照失败时，就再没有任何东西能区分“场景很难”和“链路坏了”。至于无动作探针，在 Gaia2 上根本还没有跑过。两者都不贵，两者都还没做；这就是本次交付所处位置的如实陈述。

\section{难度控制与规模化}

\subsection{难度的控制维度}

难度可以从四个方向调整。其一是信息可见性，例如主智能体能否读取角色能力说明；其二是依赖结构，从两步顺序依赖扩展到三个应用或多分支依赖；其三是通信约束，包括委派预算以及星形、链式、树形等拓扑；其四是操作规模，例如分页数量、实体数量和状态修改数量。

这些因素应逐层加入。一个有效的难度调节项需要同时满足两点：参考解仍能稳定通过，模型表现相较上一层又出现可重复的变化。若某种拓扑只是使任务无法执行，它带来的是环境故障，而不是更高难度。当前链式配置正因不满足可解性而未进入交付（\S4.2）。

\subsection{低成本环节：静态分析与批量渲染}

造题流程的前半段由几条命令完成，均不调用模型。\texttt{gaia2\_fetch} 拉取一个公开划分；\texttt{gaia2\_admission\_probe} 读取每个场景的黄金写操作与应用初始状态，推导角色集合，并执行保留信息缝场景的状态溯源判定；\texttt{gaia2\_render\_probe} 把入选场景与约束配置写成实例规格，\texttt{forge.gaia2.cli render} 再写出 Harbor 任务目录。这些工作只涉及静态分析和文件生成，在普通笔记本上即可在数分钟内完成。

这一步成本较低，是因为 \S3 中的 $W$ 只改变环境，指令 $x$ 和判定器 $J$ 可以直接复用。在读取的两个划分上，37 个信息缝场景与四实例网格组合后，一条命令即可生成 \textbf{148 个实例}。

困难在于确认这 148 个实例中有多少真正可解、可训。主要成本出现在渲染之后——而在本任务库上，其中第一项完全换了形态。

\subsection{主要成本之一：参考解换了位置，但没有消失}

若任务库的参考信息是\emph{代码}，这一节讨论的就是编写成本：Harbor 任务需要在 \texttt{solution/} 中提供参考解以证明可解，而这些路径只能一条条手写，手写参考解的数量因此成为可交付目录数的硬上限。

Gaia2 直接消除了这项成本，这也是选择该任务库最有力的实际理由。它的参考信息是一组黄金写操作，划分中的每个场景都自带，因此无需为任何一道题撰写内容就能知道“做完”意味着什么：全部 37 个场景的验收目标随数据免费获得。

它没有提供的是一条\emph{可运行}的参考路径。黄金写操作描述的是终点，而不是一条可以让智能体照着走的路线——而 Harbor 的 \texttt{solution/} 是要能执行的。所以成本没有消失，只是从“为每道题写一条参考路径”变成了“没有廉价手段证明某个场景可达”。\S4.3 指出了后果，\S7 支付了账单：由于没有参考解可跑，无约束对照成了唯一的可解性证据，而本次实验的对照在 37 个场景中失败了 35 个。于是最核心的问题——这些场景是真的难，还是这个运行时坏了——在形式上仍然悬而未决。

有两件事可以了结它，且都不昂贵。一个\textbf{黄金回放}智能体直接发出记录在案的写操作，完全不需要模型调用就能确立上界，并且在问题出自链路而非场景时会明确报错。一个\textbf{无动作探针}——本项目在别处已经在跑的下界测量——则给出区间的另一端。两者合起来，就能把 \S7 中的每个数字从“具有启发性”变成“可以解读”。它们尚未被执行，这是本次交付最大的缺口。

\subsection{主要成本之二：多随机种子验证}

配置是否有效取决于多次运行的结果分布，不能由单次结果决定（\S4.4）。按当前设定，每个实例至少需要 5 次重复。37 个场景、4 个实例和 5 个随机种子需要 740 次运行，而本文报告的实验只跑了 \textbf{148} 次——单种子，证据文件在自己的说明字段里就是这么写的。

实际瓶颈不是算力，而是另外两件事。第一，上限来自 API 速率而非计算：实验以 3 路并发运行，因为组织的每分钟令牌配额使更高并发适得其反，所以墙钟时间由速率限制决定，而不是由工作量决定。第二，验证成本会随配置数量线性增长：每增加一个调节项，都需要在全部场景上重新进行多随机种子对照。这正是“先测良率再扩规模”的理由，而 \S7 展示了从反方向违反它是什么样子：在还没有人确认对照能够完成任何一个场景之前，148 次运行就已经花在了四个旋钮上。每个场景多出的那几个实例本身没有错，错的是\emph{顺序}。

\subsection{规模上限：37 个场景，但只有 12 种信息缝结构}

即使补齐上述两项成本，任务库本身的多样性仍然有限，而且明显低于 37。

这 37 个信息缝场景并不是 37 个独立样本。按信息缝签名——即溯源判定找到的全部\emph{（来源应用 $\rightarrow$ 目标应用）}边的集合——归类，它们收敛为 \textbf{12 种结构}，其中三种就覆盖了 37 个中的 25 个（表~\ref{tab:combos}）。

\begin{table}[H]
\centering
\small
\caption{37 个信息缝场景的缝签名分布}
\label{tab:combos}
\begin{tabularx}{\textwidth}{X c}
\toprule
信息缝签名 & 场景数\\
\midrule
Contacts $\rightarrow$ Cabs、InternalContacts $\rightarrow$ Cabs & 10\\
Contacts $\rightarrow$ Messages、InternalContacts $\rightarrow$ Messages & 9\\
RentAFlat $\rightarrow$ Cabs & 6\\
Emails $\rightarrow$ Calendar & 2\\
Calendar $\rightarrow$ Emails & 2\\
Contacts/InternalContacts/RentAFlat $\rightarrow$ Cabs & 2\\
其余 6 种签名，各 1 个场景 & 6\\
\midrule
合计 & \textbf{37} 个场景，\textbf{12} 种结构\\
\bottomrule
\end{tabularx}
\end{table}

按端点统计会更紧。语料中的每一条缝都终止于四个应用之一——Cabs（38 条边）、Messages（20）、Calendar（6）、Emails（3）——并起始于六个应用之一，其中 Contacts 与 InternalContacts 合计 50 条且几乎总是同时出现，因为它们本就是保存同一事实的两个通讯录。直白地说，这个任务库能提供的信息缝几乎只有一种模式：\emph{某个人的标识或住址，从通讯录类应用流向预订或消息类应用。}按类别看也同样偏斜：adaptability 30 个、time 4 个、execution 3 个。

因此当前可训练的多样性更接近 12 种结构和一个主导模式，而不是 37 个场景，更不是 148 个实例。配置数量、场景数量和结构数量需要分别统计。公开的五个能力类别中有两个完全没有贡献，原因 \S4.4 已经给出：一个落在\emph{回答}而不是写操作里的事实，对溯源判定不可见。因此扩大语料只有两条路——拉取尚未读取的划分（成本很低，但目前只是外推而非测量），或者把信息缝判定扩展到可以看见读取行为的证据上。缺的不是旋钮，而是源场景和一个视野更宽的信息缝判定。

\subsection{准入条件：本方法对任务库的要求}
\label{sec:substrate}

Gaia2 是本方法的输入，而不是它的前提。方法对任务库只提出四项要求，缺少任何一项，它都会退回到“自己设计判定器”，而这正是 \S4.1 所论证的、本方法存在的理由所禁止的：

\begin{enumerate}
\setlength{\itemsep}{0.15em}
\item \textbf{不是你写的程序化判定器。} 它使奖励能够批量计算，也避免重新定义“协作是否成功”。它的全部价值恰恰在于：已经有别人替你验证过它。
\item \textbf{真实、可执行、可安装。} 而不是对某个世界的描述。一个无法 \texttt{pip install} 或 \texttt{docker pull} 的任务库，最后一定由你自己补完；补完之后它就是你设计的，第一项要求也随之作废。
\item \textbf{任务级参考信息。} 它既用于推导角色集合，也用于证明任务可解。Gaia2 完整提供了前者，后者却只提供了一半——黄金写操作描述的是终点而不是路线，\S6.3 就是这笔代价。
\item \textbf{可切开的工具边界。} 不同角色需要拥有彼此隔离的工具或数据；否则主智能体可以绕过协作，其他约束也失去意义。
\end{enumerate}

\textbf{准入是必要条件，而非充分条件。} 一个任务库可以四项全部满足，却依然产不出可用数据，因为这四项都没有度量信息缝有多深：判定器可以完美无缺，而每条任务仍然坍缩到单一角色上。这正是信息缝分析（\S3.1）必须是一次测量、而不是一个假设的原因；Gaia2 把这一点说得非常清楚——\texttt{adaptability} 划分有 99\% 是多应用的，而只有 19\% 存在信息缝。信息缝也是“换环境”这件事不再免费的地方：SWE-smith 拥有这份名单上最强的判定器，其规模化也正依赖第一项要求\cite{yang2025swesmith}，但它的缝是代码模块，因此在它上面做划分，得到的是一道多了几个智能体的软件工程题，而不是编排题。

\textbf{准入同样不度量难度}，而本次交付就是证明。Gaia2 轻松满足四项要求，它的信息缝良率也是测出来而非假设的，148 个实例却依然没有产出任何可认证的训练数据——因为这些要求完全没有涉及“一个运行时能否把智能体带着走完这么长的场景”。准入条件筛的是\emph{任务库}，其中没有任何一条筛\emph{运行链路}，而 \S7 会论证这次实验正是输在这里。

\textbf{关于普适性，把话说清楚。} 本方法中没有任何一步是为 Gaia2 定制的：准入条件是任务库的性质，角色集合读自参考信息实际触碰的工具，约束作用于访问面，而有效性判据（\S4.3）只需要一个分数、一个下界和一个对照。这是一个从构造出发的论证——读代码即可核验——而它现在有了单个案例之外的一点实证支持：两条挖掘规则被迁移到一个参考信息形态完全不同的任务库上（从参考\emph{代码}变成黄金\emph{写操作}），并在另一侧产出了可测量的信息缝分布。\textbf{尚未}迁移成功的是一次端到端的旋钮定价，原因不是方法没能迁移，而是这次实验从未确立自己的对照。因此普适性主张现在一半是测量、一半仍是构造性论证，本文标出这条界线，而不是让成功的那一半代表整体。合理的扩展顺序是：先在当前任务库上确立上界与下界（\S6.3），再把网格跑到五个种子，然后才向外扩展——而不是在没人证明能跑完的场景上继续叠加配置。

\section{数据样本与实验结果}

当前交付的样本是 \S3.1 中入选的 37 个信息缝场景，渲染为 148 个实例。表~\ref{tab:samples} 列出其中三个，用以覆盖挖掘所得的结构跨度。

\begin{table}[H]
\centering
\small
\caption{三个入选场景及其挖掘出的结构（\texttt{sweep/gaia2\_cells.json}）}
\label{tab:samples}
\begin{tabularx}{\textwidth}{>{\raggedright\arraybackslash}p{2.7cm} >{\raggedright\arraybackslash}p{2.7cm} X c c}
\toprule
场景 & 信息缝 & 角色集合 & 写操作 & 深度\\
\midrule
\texttt{..\_30\_68r6vs} & Calendar $\rightarrow$ Cabs & Cabs、Calendar、Messages & 5 & 3\\
\texttt{..\_25\_b932pi} & Contacts $\rightarrow$ Cabs & 另加 Contacts、InternalContacts & 5 & 3\\
\texttt{..\_24\_srcjx4} & Contacts $\rightarrow$ Cabs & 再加 Shopping & 14 & 4\\
\bottomrule
\end{tabularx}
\end{table}

相比本项目起步时那个单一题族，这些场景的时间跨度更长、宽度也更大：角色集合最多 7 个应用，黄金写操作最多 16 条，其中一个场景在单次扇出中就提供了 12 条可并行的写操作。其中三个已渲染为可直接运行的 Harbor 任务目录，全部 148 个则以实例规格的形式存在。

我们在每个场景上比较四个实例。\texttt{control}（\texttt{open-docs-binf}）让单个智能体持有全部工具；\texttt{context}（\texttt{star-docs-binf}）强制委派，但允许主智能体读取角色能力；\texttt{discovery}（\texttt{star-names-binf}）隐藏能力说明，只留角色名；\texttt{economy}（\texttt{star-docs-b}$N$）恢复能力文档，并加上按任务推导的软性委派目标。主智能体与专属智能体同为 \texttt{gpt-5.6-sol}，模型档位在构造上保持对等。\textbf{单种子。}

\begin{table}[H]
\centering
\small
\caption{完整实验：37 个场景 $\times$ 4 个实例 $=148$ 次运行，单种子（\texttt{sweep/gaia2\_full\_campaign.json}）}
\label{tab:results}
\begin{tabularx}{\textwidth}{>{\raggedright\arraybackslash}p{2.2cm} c c c X}
\toprule
实例 & 成功数 & 平均部分分 & 平均委派次数 & 转动的旋钮\\
\midrule
control & 2/37 & 0.251 & 0 & 无：单个智能体持有全部应用\\
discovery & 3/37 & 0.247 & 6.3 & 角色能力不可见\\
context & 2/37 & 0.253 & 6.8 & 权限划分本身\\
economy & 2/37 & 0.282 & 7.6 & 软性 \texttt{b}$N$ 委派目标\\
\bottomrule
\end{tabularx}
\end{table}

\textbf{这张表里的任何一个旋钮都无法解读，原因就在第一行。}四个均值彼此相差不超过 $0.035$，成功数处处是 37 中的 2--3 个（包括无约束组），而约束\emph{最多}的那一组均值反而最高。\S4.3 的接受判据按其第三个条件否决了这张表——$r_{\mathrm{open}}$ 没有越过下界——而这条否决本身就是本文要报告的结果。若从这些数字中得出“权限划分是免费的，委派预算还有帮助”，那不过是用更好的工具和更多的运行次数，把 \S4.1 存在的意义所要防止的错误重犯一遍。

有三项证据表明这个排序是噪声而非信号。

\textbf{其一，对照的行为不像一个上界。}在 \textbf{37 个场景中有 6 个，某个受约束组的得分高于无约束组}——持有全部工具的智能体，反而不如一个必须盲目委派的智能体。在语料中宽度最大的 \texttt{..\_24\_srcjx4} 上，对照一个动作都没有发出就投降了，得分 $0.0$，而 \texttt{discovery} 组覆盖了一半的黄金写操作。一个会输给自己盲化版本的对照，度量的不是难度。

\textbf{其二，判定器与场景总体是混淆的。}带回复条件的场景由软判定器评分，无条件的场景由脚本判定器评分，这使语料被切成了不均等的两半：

\begin{table}[H]
\centering
\small
\caption{按判定器统计的成功数——以及各自看到的、互不相交的场景集合}
\label{tab:judges}
\begin{tabularx}{\textwidth}{X c c}
\toprule
判定器 & 实例数 & 成功数\\
\midrule
\texttt{gpt-5.6-sol} 软判定器（带回复条件的场景） & 128 & \textbf{0}\\
脚本写操作检查（无条件场景） & 20 & \textbf{9}\\
\bottomrule
\end{tabularx}
\end{table}

128 比 0、20 比 9，差距很大，也很容易被读成“软判定器过严”。但它不能这样读：两个判定器从未评过同一批场景，因此这个差距同时为判定器和场景总体定价，而不为其中任何一个单独定价。要把它们分开只需一次廉价实验——让两个判定器评同一批轨迹——而它还没有做。在此之前，上面每一个按实例统计的均值都混合了两个不同的总体。

\textbf{其三，也是最有用的一点：轨迹在协作受审之前就已经死了。}\code{gaia2_credit.py} 为每条轨迹标注首个故障点，过程确定且不调用模型：

\begin{table}[H]
\centering
\small
\caption{148 次运行首次出错的位置（\texttt{sweep/gaia2\_credit.json}）}
\label{tab:pivots}
\begin{tabularx}{\textwidth}{>{\raggedright\arraybackslash}p{3.0cm} c X}
\toprule
首个故障 & 运行数 & 含义\\
\midrule
\texttt{wrong-arguments} & 45 & 调用发出了，但参数与黄金不符\\
\texttt{malformed} & 40 & 模型的工具调用语法从未解析成功\\
\texttt{environment-stop} & 33 & 环境先一步结束了这一幕\\
\texttt{refusal} & 11 & 专属智能体拒绝了请求\\
\texttt{surrender} & 9 & 主智能体使用了 FAIL 动作\\
\texttt{false-outage} & 7 & 专属智能体谎称某个工具坏了\\
\bottomrule
\end{tabularx}
\end{table}

\textbf{格式错误与环境结束合计 148 次中的 73 次——49\%。}两者都不是协作失败。格式错误是运行时没有强制、或模型没有遵守的语法约定；环境结束则是场景的时钟走完了。只看对照组，情况更为严峻：37 次对照运行中有 19 次终止于 \texttt{environment-stop}、13 次终止于 \texttt{malformed}，也就是说 \textbf{37 次无约束运行里有 32 次根本没有走到任务可以被完成的状态}。这是一个披着难度外衣的链路结果，而这正是 \S4.3 第三个条件所要拦截的混淆。

因此如实的总结很短。挖掘规则完成了迁移并产出了可测量的信息缝分布；运行时端到端跑完了 148 个实例并记录了完整轨迹；而接受判据作用在这些实例上时，\emph{拒绝给出判定}——这是正确的。这次实验定价的不是约束，而是运行时，代价是 49\% 的轨迹。

\section{讨论与后续延展}

\subsection{当前结果与下一步验证}

下一步的顺序直接来自 \S7，而且没有一项昂贵的选择排在前面。

\textbf{先确立上界与下界。}一个直接发出记录在案的写操作的黄金回放智能体，不需要任何模型调用，就能回答实验无法回答的问题：一个对照失败了的场景，在这个运行时里究竟是否可达？而无动作探针——本项目在别处已经在跑的下界测量——提供区间的另一端。两者在 Gaia2 上都还没有跑过；在跑完之前，任何实例都无法认证，任何旋钮都无法定价。

\textbf{先修运行时，再往上加东西。}40 次格式错误和 33 次环境结束占了语料的 49\%，而它们属于可解决的工程问题而非研究问题：调用语法可以强制而不是请求，时钟预算可以按每个场景的时间跨度去度量。在此之前投在第五个旋钮上的每一次运行，都只是在测量评测链路。

\textbf{把判定器和场景总体分开。}让两个判定器评同一批轨迹只需一次遍历，就能把表~\ref{tab:judges} 从一个无法解读的差距，变成一次对判定器严格程度的测量。

\textbf{最后才是种子。}37 个场景、4 个实例、5 次重复是 740 次运行；这笔钱值得花，但要等到对照越过下界之后，而不是之前。

还有第二个性质不同的缺口。\textbf{本文把方法写成与任务库无关}，而本次交付是对这一主张的第一次真实检验，因为 Gaia2 的参考信息形态与参考代码完全不同。挖掘规则迁移成功了：角色推导与信息缝判定都在另一侧产出了可测量的分布，这是实打实的证据，也是本文最强的结果。没有迁移成功的是一次端到端的旋钮定价——原因不是方法未能迁移，而是这次实验从未确立自己的对照。因此普适性主张现在一半是测量、一半仍是构造性论证，本文标出这条界线，而不是让成功的那一半代表整体。

终态奖励可以判断任务是否完成，却不能区分失败来自拆解、角色选择还是信息表达。表~\ref{tab:pivots} 的故障点标注是走向归因的第一步——它确定、不含模型，并且已经能够区分“从未发出可解析的调用”与“调用了正确工具但参数有误”。但它还不是能力归因；在其准确率与人工标注对照测量之前，它不应进入奖励——这与本文对任何自行设计的判定器所用的规则相同。

\subsection{从空间约束延展到时间约束}

当前工作主要改变协作的空间结构：工具属于谁、信息对谁可见、消息可以沿哪些路径传递。下一阶段可以加入时间维度，不再在任务开始时一次给出全部信息，而是在智能体主动询问或环境达到特定状态后再揭示。这样可以训练隐藏需求识别和用户干预下的重规划，同时继续复用现有终态验证。

最小实现可以从同一任务生成一组成对样本。对照组提供完整指令；实验组只给出符合真实使用习惯的不完整请求，把联系人范围、日期或兼容偏好保存在辅助容器的隐藏用户状态中。主智能体需要先识别信息缺口，再通过询问获得这些约束。两组任务共享初始环境和最终判定器，因此不必直接评价“问题问得好不好”：遗漏关键问题会导致终态错误，而无差别追问则会消耗有限的交互预算。

再进一步，可以在任务执行过程中加入用户干预。例如，智能体完成部分工作后，用户补充新的范围或限制，主智能体需要更新计划、复用仍然有效的结果，并把变化同步给相关角色。这里要区分两类信息：若后续补充只是初始需求中未说清的部分，原有判定器可以保持不变；若用户真正改变了目标，就必须为新终态提供同样可靠的环境验证，不能只根据对话判断是否适应成功。

这条时间轴扩展可以表示为
\[
  \tau_{c,h}=(x_0,\mathcal{E}_c,U_h,\Gamma_h,J),
\]
其中 $U_h$ 表示只存在于辅助容器的隐藏用户状态，$\Gamma_h$ 描述信息被询问、揭示或修改的时机。这里仍需保持现有隔离边界：隐藏状态不能进入主智能体的文件系统，交互记录不能由智能体自行伪造，最终奖励继续由可重放的环境状态决定。只有当完整指令、隐藏需求和中途干预三种配置形成稳定差异后，才适合把时间维度纳入正式课程。

\section{结语}

这项工作的核心，是从已有真实任务中识别跨角色依赖，再通过权限和信息划分让这些依赖成为完成任务的必要条件。主智能体负责规划、委派和传递上下文，专属智能体只执行局部工作；最终判定仍由 Gaia2 自己的写操作判定器给出。

本文确立了两件事。挖掘规则\textbf{迁移到了一个参考信息形态完全不同的任务库}上——从参考代码变为黄金写操作——并在两个划分上产出了可测量的信息缝分布：160 中的 17 个与 160 中的 30 个，而对应的多应用率分别是 69\% 与 99\%。整条链路也确实端到端跑通了：挖掘 37 个场景，渲染 148 个实例，执行 148 次运行并保留完整轨迹，其中三个打包为可直接运行的 Harbor 任务目录。

有一件事没有确立。\textbf{这些约束能否形成稳定的难度梯度，仍未得到证明，而本次实验也没有让它更接近证明。}四个实例彼此相差不超过 $0.035$，无约束对照在 37 个场景中只成功 2 个，49\% 的运行终止于格式错误的调用或环境提前结束。接受判据拒绝认证任何一个实例——这正是判据在正常工作：它被设计成在对照没有越过下界时保留判断，而它做到了。

因此下一步不是扩大规模，而是补上这次实验所缺的两个参照——黄金回放给出的上界与无动作探针给出的下界——并修复运行时中正在把轨迹输给语法和时钟的那一半。只有在此之后，种子、更多划分和新旋钮才买得到东西。一个连自己的对照都还跑不完的造题链路，没有理由去增加第五种约束；而本次交付产出的最有价值的东西，正是那份说明了这一点的证据。

\begingroup
\scriptsize
\begin{thebibliography}{9}
\setlength{\itemsep}{0.05em}
\setlength{\parskip}{0pt}

\bibitem{meta2025are}
Meta AI.
\newblock ARE: Scaling Up Agent Environments and Evaluations.
\newblock arXiv:2509.17158, 2025. 提出 Gaia2 基准；软件包名为 \texttt{meta-agents-research-environments}。
\newblock \url{https://arxiv.org/abs/2509.17158}

\bibitem{yang2025swesmith}
John Yang, Kilian Lieret, Carlos E. Jimenez, et al.
\newblock SWE-smith: Scaling Data for Software Engineering Agents.
\newblock arXiv:2504.21798, 2025.
\newblock \url{https://arxiv.org/abs/2504.21798}

\bibitem{zhu2025multiagentbench}
Kunlun Zhu, Hongyi Du, Zhaochen Hong, et al.
\newblock MultiAgentBench: Evaluating the Collaboration and Competition of LLM Agents.
\newblock \emph{ACL 2025}, arXiv:2503.01935.
\newblock \url{https://arxiv.org/abs/2503.01935}

\bibitem{cemri2025mast}
Mert Cemri, Melissa Z. Pan, Shuyi Yang, et al.
\newblock Why Do Multi-Agent LLM Systems Fail?
\newblock arXiv:2503.13657, 2025.
\newblock \url{https://arxiv.org/abs/2503.13657}

\bibitem{harbor2026tasks}
Harbor Framework.
\newblock Task Structure: Creating and Running Tasks for Agentic Environments.
\newblock Documentation, accessed 2026-07-16.
\newblock \url{https://www.harborframework.com/docs/tasks}

\end{thebibliography}
\endgroup

\end{document}
